\documentclass[12pt,oneside,a4paper]{article}
\usepackage[english]{babel}
\usepackage[utf8x]{vietnam}
% package for including graphics with figure-environment
\usepackage{graphicx}
\usepackage{amsmath,amssymb,exscale,eucal,amsthm, mathrsfs}

%\DeclareMathSizes{12}{15}{9}{9}
\usepackage[table,xcdraw]{xcolor}
\usepackage{tabularx}

\usepackage[top=3.5cm, bottom=3cm, left=2cm, right=2cm] {geometry}
% colors for hyperlinks
% colored borders (false) colored text (true)
\usepackage{wrapfig}
\usepackage{color}
%\usepackage{caption}
% package for bibliography
%\usepackage[round, sort, numbers]{natbib}
% package for header
\usepackage[automark]{scrpage2}
\pagestyle{scrheadings}
\ihead[]{Tìm hiểu về SNMP và công cụ giám sát mạng}
\ohead[]{\today}
\cfoot[]{\pagemark} 
\setheadsepline[122mm]{0.3mm}

\usepackage{enumitem}


\begin{document}
		
		\addtolength{\leftskip}{5em} \addtolength{\rightskip}{2em}
		
		\begin{wrapfigure}[5]{l}{5em}\centering
			\includegraphics[width=0.1\textwidth]{bkhn.jpg}
		\end{wrapfigure}
		
		\noindent  \\
		 \\
		 \\
		Trường Đại học Bách Khoa Hà Nội \\
		Viện Toán ứng dụng và Tin học
		
		\addtolength{\leftskip}{-5em} \addtolength{\rightskip}{-2em}% Restore margins

		\title{\Huge{Báo cáo bài tập lớn}\\ \large{Môn: Thiết kế, cài đặt và quản trị mạng máy tính}\\}
	
	
	
	\vspace{1cm}
	
	\author{Đề tài: Tìm hiểu về SNMP và công cụ giám sát mạng Cacti}
	
	\vspace{1cm}
	
	% name of the course and module
	\date{
	\vspace{1cm}
	\begin{table}[h]
		\centering
		\begin{tabular}{rl}
			\large Sinh viên thực hiện: & \large Đào Tuấn Anh (MSSV: 20121178)\\
                    						& \large Vũ Thanh Tùng (MSSV: 20124917)
		\end{tabular}
	\end{table}	
	\vspace{1cm}
	\begin{table}[h]
		\centering
		\begin{tabular}{rl}
			\large Giáo viên hướng dẫn: & \large ThS. Nguyễn Tuấn Dũng\\
		\end{tabular}
	\end{table}
	\vspace{1cm}
	Hà Nội, \today
	}
	\global\let\newpagegood\newpage
	\global\let\newpage\relax\maketitle
	\global\let\newpage\newpagegood
	
	\setlength{\parindent}{0pt}

	\newpage
	\tableofcontents
	\newpage
	

\section*{Lời mở đầu}
\addcontentsline{toc}{section}{\protect\numberline{}Lời mở đầu}%
Sự phát triển không ngừng nghỉ của công nghệ thông tin đã tạo ra những thay đổi lớn trong cơ sở hạ tầng, lực lượng sản xuất, tính chất lao động, cấu trúc kinh tế và cả cách thức quản lý trong các lĩnh vực của xã hội. Một trong những yếu tố quyết định, cũng là một nhu cầu tất yếu là việc trao đổi thông tin giữa các thiết bị thông qua mạng máy tính. Sự tăng nhanh về kích cỡ của các hệ thống mạng đã và đang đưa ra thách thức không nhỏ cho người quản trị trong việc theo dõi và quản lý.\\

Do đó, các hệ thống giám sát an toàn mạng ra đời, đóng vai trò quan trọng, không thể thiếu trong hạ tầng công nghệ thông tin của các cơ quan, đơn vị, tổ chức. Hệ thống này cho phép thu thập, chuẩn hóa, lưu trữ và phân tích tương quan toàn bộ các sự kiện được sinh ra trong hệ thống mạng của tổ chức. Từ đó người quản trị có thể dễ dàng đưa ra phân tích, thống kê, cảnh báo, nắm bắt thực trạng tài nguyên, khắc phục sự cố, đồng thời ngăn chặn sớm các mối nguy hiểm cho hệ thống.\\

Nội dung dưới đây có đề cập đến giao thức quản lý mạng SNMP (Simple Network Management Protocol). SNMP là một tập hợp các giao thức không chỉ cho phép kiểm tra các thiết bị mạng như router, switch hay server có đang vận hành mà còn hỗ trợ vận hành các thiết bị này một cách tối ưu, ngoài ra SNMP còn cho phép quản lý các thiết bị mạng từ xa. Bên cạnh đó, báo cáo còn đề cập đến Cacti, một hệ thống quản lý tài nguyên mạng mã nguồn mở. Phần mềm này đáp ứng nhu cầu quản lý mạng một cách toàn diện với nhiều tính năng linh hoạt vượt trội.\\

Em xin chân thành cảm ơn thầy Nguyễn Tuấn Dũng đã tận tình truyền đạt kiến thức trong suốt môn học. Tuy nhiên trong khoảng thời gian cho phép, việc nghiên cứu và trình bày không thể tránh khỏi nhầm lẫn thiếu sót, em mong nhận được sự giúp đỡ đóng góp của thầy cô và bạn bè để bài báo cáo được hoàn thiện hơn.\\
\vspace{1cm}
\textit{Sinh viên}\\

Đào Tuấn Anh - Vũ Thanh Tùng
\newpage

\section{Tìm hiểu về giao thức SNMP}
\subsection{Tổng quan về quản lý hệ thống mạng}
\subsubsection{Khái niệm quản lý hệ thống mạng}
Quản lý hệ thống mạng là việc thực hiện tập hợp các chức năng cần thiết cho việc kiểm soát, lập kế hoạch, phân bổ, triển khai, phối hợp và giám sát các nguồn tài nguyên trong hệ thống. Nó thường được áp dụng cho các mạng quy mô lớn như mạng máy tính và mạng viễn thông, đề cập đến việc duy trì và quản trị mạng ở cấp cao nhất.
\subsubsection{Khái niệm giám sát mạng}
Giám sát mạng là việc sử dụng một hệ thống liên tục theo dõi hoạt động của một mạng máy tính, thông báo ngay lập tức và cho phép người quản trị mạng biết trong trường hợp xảy ra sự cố ảnh hưởng đến tốc độ hoặc lỗi mạng.\\

Giám sát mạng là một phần trong quản lý hệ thống mạng.\\

Một hệ thống giám sát mạng sẽ theo dõi hệ thống mạng gây ra bởi quá tải hay máy chủ bị sập, bởi các kết nối trong mạng hoặc từ các thiết bị khác. Các thông số đánh giá thường gặp là thời gian phản hồi, hiệu lực và thời gian vận hành tối thiểu trung bình (uptime). Ngoài ra còn có các thông số phổ biến khác như độ ổn định hay độ tin cậy.\\

Ví dụ, để kiểm tra một máy chủ Web, phần mềm giám sát sẽ gửi đều đặn các truy vấn HTTP để lấy nội dung trang web. Ở các máy chủ thư, các gói tin kiểm tra sẽ được gửi đều đặn thông qua SMTP, và nhận lại qua IMAP hoặc POP3.
\subsection{Giao thức quản lý mạng SNMP (Simple Network Management Protocol)}
\subsubsection{Giới thiệu giao thức SNMP}
\subsubsection{Ưu điểm của giao thức SNMP}
\subsubsection{Các phiên bản giao thức SNMP}
\subsubsection{Mô hình hệ thống mạng sử dụng SNMP}
\subsubsection{Kiến trúc giao thức SNMP}
\subsubsection{Các cơ chế bảo mật trong SNMP}

\section{Công cụ giám sát mạng Cacti}
\subsection{Giới thiệu về công cụ giám sát mạng Cacti}
\subsection{Kiến trúc phần mềm Cacti}
\subsection{Cài đặt Cacti trên Windows}
\subsection{Triển khai và ứng dụng Cacti trong giám sát mạng}
\subsubsection{Mô hình triển khai}
\subsubsection{Cấu hình Cacti}

\section*{Kết luận}
\addcontentsline{toc}{section}{\protect\numberline{}Kết luận}%

\newpage
\section*{Tài liệu tham khảo}
\addcontentsline{toc}{section}{\protect\numberline{}Tài liệu tham khảo}%
%\bibliographystyle{ksfh_nat}
%	\bibliography{references} % expects file "references.bib"
\begin{enumerate}
\item{Wikipedia, "Simple Network Management Protocol", https://en.wikipedia.org/wiki/SNMP}
\item{Wikipedia, "Network monitoring", https://en.wikipedia.org/wiki/Network\_monitoring}
\item{TBT - VNCERT, "Tổng quan về hệ thống giám sát mạng", http://www.vncert.gov.vn/}
\item{Diệp Thanh Nguyên, "SNMP Toàn tập", https://sites.google.com/site/snmptoantap/home}
\end{enumerate}
	

\end{document}